\documentclass[11pt,a4paper]{article}
% main (standalone preamble for editing)
% vim: nowrap: spell:
% ══════════════════════════════════════════════════════════════════════════════
% Speed Maths --- shared preamble
% Included by every sheet and answer file via:  % main (standalone preamble for editing)
% vim: nowrap: spell:
% ══════════════════════════════════════════════════════════════════════════════
% Speed Maths --- shared preamble
% Included by every sheet and answer file via:  % main (standalone preamble for editing)
% vim: nowrap: spell:
% ══════════════════════════════════════════════════════════════════════════════
% Speed Maths --- shared preamble
% Included by every sheet and answer file via:  \input{../../shared/preamble}
% Editing THIS file changes the look of every sheet/answer across every pillar.
% ══════════════════════════════════════════════════════════════════════════════

\usepackage[a4paper, top=25mm, bottom=25mm, left=22mm, right=22mm]{geometry}
\usepackage{amsmath, amssymb}
\usepackage{enumitem}
\usepackage{booktabs}
\usepackage{array}
\usepackage{parskip}
\usepackage{microtype}
\usepackage{fancyhdr}
\usepackage{titlesec}
\usepackage{mdframed}
\usepackage{xcolor}
\usepackage[hidelinks]{hyperref}
\usepackage{graphicx}
\usepackage{tikz}
\usepackage{pgfplots}
\pgfplotsset{compat=1.18}

% ── Colours ───────────────────────────────────────────────────────────────────
\definecolor{darkgray}{gray}{0.15}
\definecolor{medgray}{gray}{0.45}
\definecolor{lightgray}{gray}{0.93}
\definecolor{boxgray}{gray}{0.88}

% ── Section formatting ──────────────────────────────────────────────────────────
\titleformat{\section}[block]
  {\normalfont\large\bfseries}{}
  {0pt}{}[\vspace{2pt}\hrule\vspace{6pt}]
\titlespacing*{\section}{0pt}{14pt}{4pt}

% ── List formatting ─────────────────────────────────────────────────────────────
\setlist[enumerate,1]{label=\textbf{\arabic*.}, leftmargin=2em, itemsep=10pt}

% ── Branding constants (edit here, not per-file) ────────────────────────────────
\newcommand{\SpeedExamLine}{\textbf{Speed Maths:} Competition \& Exam Prep}
\newcommand{\SpeedCredit}{Series by \href{https://www.linkedin.com/in/cerealdev/}{David Akinyele-Aje}}

% ── Header / footer ──────────────────────────────────────────────────────────────
% Usage (once, after \input{../../shared/preamble} and before \begin{document}):
%   \SpeedHeader{Algebra}{3}
% First arg  = pillar name shown as "Speed <Pillar> --- Daily Drill"
% Second arg = worksheet number
\pagestyle{fancy}
\newcommand{\SpeedHeader}[2]{%
  \fancyhf{}%
  \renewcommand{\headrulewidth}{0.4pt}%
  \setlength{\headheight}{14pt}%
  \fancyhead[L]{\small\textbf{Speed #1 --- Daily Drill}}%
  \fancyhead[R]{\small Worksheet \##2}%
  \fancyfoot[L]{\small \SpeedExamLine}%
  \fancyfoot[C]{\small\thepage}%
  \fancyfoot[R]{\small \SpeedCredit}%
  \renewcommand{\footrulewidth}{0.4pt}%
}

% ── Title block (used at the top of every sheet AND answer file) ───────────────
% Usage: \SpeedTitleBlock{Daily Algebra Drill \#3}{Jane Doe}
% %% AGENTS: When generating a new sheet, ask the human contributor for the name they want credited in argument 2.
% For answer files, pass the "--- Answers \& Investigations" suffix in the arg.
\newcommand{\SpeedTitleBlock}[2]{%
  \begin{center}
    {\LARGE\bfseries #1}\\[4pt]
    {\normalsize\itshape Sheet by #2}\\[6pt]
    \hrule\vspace{4pt}
    \begin{tabular}{llcc}
      \toprule
      \textbf{Section} & \textbf{Topic} & \textbf{Questions} & \textbf{Target time}\\
      \midrule
      A & Rapid Recognition          & 10 & 2 min 30 sec \\
      B & Manipulation Drills        & 10 & 8 min        \\
      C & Substitution \& Structure  & 8  & 10 min       \\
      D & Challenge                  & 5  & 15 min       \\
      \bottomrule
    \end{tabular}
    \vspace{4pt}
    \hrule
  \end{center}
}

% ── Sheet metadata: topic + the day's new tools ────────────────────────────────
% Usage (immediately after \SpeedTitleBlock, sheets only):
%   \SpeedMeta{Expanding, factorising, and the identity toolkit}
%             {difference of squares, sum/product form, completing the square}
%
% Arg 1 = the sheet's topic in one line. The website lists this instead of
%         "Sheet 03", so write the thing a reader would actually search for.
% Arg 2 = comma-separated, at most five: the tools this sheet introduces that
%         earlier sheets in the pillar did not. These become the website's tags.
%         Leave out anything the reader already met — the tags are meant to say
%         what is new today, not to summarise the sheet.
%
% Both are parsed straight out of this file by tools/build_website.py, so the
% sheet stays the single source of truth and the site cannot drift from it.
%
% This renders NOTHING. It is a declaration for the build script, not a piece of
% the sheet's layout: adding it to a sheet must not change that sheet's PDF, so
% the 25 existing sheets can be annotated without a single one of them being
% reissued. If the toolkit should also appear on the page, that is a separate
% decision about the sheet design — make it deliberately, here, not by leaving
% this macro printing something nobody asked it to.
\newcommand{\SpeedMeta}[2]{}

% ── Answer / Method / Investigate-further boxes (answer files only) ────────────
\newmdenv[
  backgroundcolor=lightgray,
  linecolor=medgray,
  linewidth=0.5pt,
  leftmargin=0pt, rightmargin=0pt,
  innerleftmargin=8pt, innerrightmargin=8pt,
  innertopmargin=5pt, innerbottommargin=5pt,
  skipabove=4pt, skipbelow=4pt
]{ansbox}

\newmdenv[
  backgroundcolor=white,
  linecolor=darkgray,
  linewidth=0.8pt,
  leftline=true, rightline=false, topline=false, bottomline=false,
  leftmargin=0pt, rightmargin=0pt,
  innerleftmargin=8pt, innerrightmargin=6pt,
  innertopmargin=4pt, innerbottommargin=4pt,
  skipabove=4pt, skipbelow=6pt
]{invbox}

\newcommand{\ans}[1]{%
  \begin{ansbox}
    \textbf{Answer:} #1
  \end{ansbox}}

\newcommand{\method}[1]{%
  {\small\textbf{Method:} #1}\par}

\newcommand{\inv}[1]{%
  \begin{invbox}
    \textbf{\small Investigate further:} {\small #1}
  \end{invbox}}

% ── Misc helpers ────────────────────────────────────────────────────────────────
\newcommand{\qn}[1]{\textbf{#1.}\;}

% Mistake-linking: attach a Top 5 Patterns / Common Traps bullet to the question
% label(s) in this sheet where it actually shows up, e.g. \seealso{B6, D2}
\newcommand{\seealso}[1]{\hfill\textit{\footnotesize(see #1)}}

% ── Graph options macro ─────────────────────────────────────────────────────────
% Usage: \GraphOptions{tikz code for A}{tikz code for B}{tikz code for C}{tikz code for D}
\newcommand{\GraphOptions}[4]{%
  \begin{center}
    \begin{tabular}{cc}
      \textbf{A)} & \textbf{B)} \\
      #1 & #2 \\[10pt]
      \textbf{C)} & \textbf{D)} \\
      #3 & #4
    \end{tabular}
  \end{center}
}

% ── Closing motivational rule + cross-promo footer (sheets only) ────────────────
% Usage: \SpeedClosing{"Quote text."}
\newcommand{\SpeedClosing}[1]{%
  \vspace{20pt}
  \begin{center}
  \rule{0.6\textwidth}{0.4pt}\\[4pt]
  {\small\itshape #1}\\[4pt]
  \rule{0.6\textwidth}{0.4pt}
  \end{center}
  \begin{center}
  {\small Found a mistake or want to contribute a sheet?}\\
  {\small Visit the open-source project at \href{https://github.com/The-CerealDev/speed-maths}{github.com/The-CerealDev/speed-maths}}
  \end{center}
}

% Editing THIS file changes the look of every sheet/answer across every pillar.
% ══════════════════════════════════════════════════════════════════════════════

\usepackage[a4paper, top=25mm, bottom=25mm, left=22mm, right=22mm]{geometry}
\usepackage{amsmath, amssymb}
\usepackage{enumitem}
\usepackage{booktabs}
\usepackage{array}
\usepackage{parskip}
\usepackage{microtype}
\usepackage{fancyhdr}
\usepackage{titlesec}
\usepackage{mdframed}
\usepackage{xcolor}
\usepackage[hidelinks]{hyperref}
\usepackage{graphicx}
\usepackage{tikz}
\usepackage{pgfplots}
\pgfplotsset{compat=1.18}

% ── Colours ───────────────────────────────────────────────────────────────────
\definecolor{darkgray}{gray}{0.15}
\definecolor{medgray}{gray}{0.45}
\definecolor{lightgray}{gray}{0.93}
\definecolor{boxgray}{gray}{0.88}

% ── Section formatting ──────────────────────────────────────────────────────────
\titleformat{\section}[block]
  {\normalfont\large\bfseries}{}
  {0pt}{}[\vspace{2pt}\hrule\vspace{6pt}]
\titlespacing*{\section}{0pt}{14pt}{4pt}

% ── List formatting ─────────────────────────────────────────────────────────────
\setlist[enumerate,1]{label=\textbf{\arabic*.}, leftmargin=2em, itemsep=10pt}

% ── Branding constants (edit here, not per-file) ────────────────────────────────
\newcommand{\SpeedExamLine}{\textbf{Speed Maths:} Competition \& Exam Prep}
\newcommand{\SpeedCredit}{Series by \href{https://www.linkedin.com/in/cerealdev/}{David Akinyele-Aje}}

% ── Header / footer ──────────────────────────────────────────────────────────────
% Usage (once, after % main (standalone preamble for editing)
% vim: nowrap: spell:
% ══════════════════════════════════════════════════════════════════════════════
% Speed Maths --- shared preamble
% Included by every sheet and answer file via:  \input{../../shared/preamble}
% Editing THIS file changes the look of every sheet/answer across every pillar.
% ══════════════════════════════════════════════════════════════════════════════

\usepackage[a4paper, top=25mm, bottom=25mm, left=22mm, right=22mm]{geometry}
\usepackage{amsmath, amssymb}
\usepackage{enumitem}
\usepackage{booktabs}
\usepackage{array}
\usepackage{parskip}
\usepackage{microtype}
\usepackage{fancyhdr}
\usepackage{titlesec}
\usepackage{mdframed}
\usepackage{xcolor}
\usepackage[hidelinks]{hyperref}
\usepackage{graphicx}
\usepackage{tikz}
\usepackage{pgfplots}
\pgfplotsset{compat=1.18}

% ── Colours ───────────────────────────────────────────────────────────────────
\definecolor{darkgray}{gray}{0.15}
\definecolor{medgray}{gray}{0.45}
\definecolor{lightgray}{gray}{0.93}
\definecolor{boxgray}{gray}{0.88}

% ── Section formatting ──────────────────────────────────────────────────────────
\titleformat{\section}[block]
  {\normalfont\large\bfseries}{}
  {0pt}{}[\vspace{2pt}\hrule\vspace{6pt}]
\titlespacing*{\section}{0pt}{14pt}{4pt}

% ── List formatting ─────────────────────────────────────────────────────────────
\setlist[enumerate,1]{label=\textbf{\arabic*.}, leftmargin=2em, itemsep=10pt}

% ── Branding constants (edit here, not per-file) ────────────────────────────────
\newcommand{\SpeedExamLine}{\textbf{Speed Maths:} Competition \& Exam Prep}
\newcommand{\SpeedCredit}{Series by \href{https://www.linkedin.com/in/cerealdev/}{David Akinyele-Aje}}

% ── Header / footer ──────────────────────────────────────────────────────────────
% Usage (once, after \input{../../shared/preamble} and before \begin{document}):
%   \SpeedHeader{Algebra}{3}
% First arg  = pillar name shown as "Speed <Pillar> --- Daily Drill"
% Second arg = worksheet number
\pagestyle{fancy}
\newcommand{\SpeedHeader}[2]{%
  \fancyhf{}%
  \renewcommand{\headrulewidth}{0.4pt}%
  \setlength{\headheight}{14pt}%
  \fancyhead[L]{\small\textbf{Speed #1 --- Daily Drill}}%
  \fancyhead[R]{\small Worksheet \##2}%
  \fancyfoot[L]{\small \SpeedExamLine}%
  \fancyfoot[C]{\small\thepage}%
  \fancyfoot[R]{\small \SpeedCredit}%
  \renewcommand{\footrulewidth}{0.4pt}%
}

% ── Title block (used at the top of every sheet AND answer file) ───────────────
% Usage: \SpeedTitleBlock{Daily Algebra Drill \#3}{Jane Doe}
% %% AGENTS: When generating a new sheet, ask the human contributor for the name they want credited in argument 2.
% For answer files, pass the "--- Answers \& Investigations" suffix in the arg.
\newcommand{\SpeedTitleBlock}[2]{%
  \begin{center}
    {\LARGE\bfseries #1}\\[4pt]
    {\normalsize\itshape Sheet by #2}\\[6pt]
    \hrule\vspace{4pt}
    \begin{tabular}{llcc}
      \toprule
      \textbf{Section} & \textbf{Topic} & \textbf{Questions} & \textbf{Target time}\\
      \midrule
      A & Rapid Recognition          & 10 & 2 min 30 sec \\
      B & Manipulation Drills        & 10 & 8 min        \\
      C & Substitution \& Structure  & 8  & 10 min       \\
      D & Challenge                  & 5  & 15 min       \\
      \bottomrule
    \end{tabular}
    \vspace{4pt}
    \hrule
  \end{center}
}

% ── Sheet metadata: topic + the day's new tools ────────────────────────────────
% Usage (immediately after \SpeedTitleBlock, sheets only):
%   \SpeedMeta{Expanding, factorising, and the identity toolkit}
%             {difference of squares, sum/product form, completing the square}
%
% Arg 1 = the sheet's topic in one line. The website lists this instead of
%         "Sheet 03", so write the thing a reader would actually search for.
% Arg 2 = comma-separated, at most five: the tools this sheet introduces that
%         earlier sheets in the pillar did not. These become the website's tags.
%         Leave out anything the reader already met — the tags are meant to say
%         what is new today, not to summarise the sheet.
%
% Both are parsed straight out of this file by tools/build_website.py, so the
% sheet stays the single source of truth and the site cannot drift from it.
%
% This renders NOTHING. It is a declaration for the build script, not a piece of
% the sheet's layout: adding it to a sheet must not change that sheet's PDF, so
% the 25 existing sheets can be annotated without a single one of them being
% reissued. If the toolkit should also appear on the page, that is a separate
% decision about the sheet design — make it deliberately, here, not by leaving
% this macro printing something nobody asked it to.
\newcommand{\SpeedMeta}[2]{}

% ── Answer / Method / Investigate-further boxes (answer files only) ────────────
\newmdenv[
  backgroundcolor=lightgray,
  linecolor=medgray,
  linewidth=0.5pt,
  leftmargin=0pt, rightmargin=0pt,
  innerleftmargin=8pt, innerrightmargin=8pt,
  innertopmargin=5pt, innerbottommargin=5pt,
  skipabove=4pt, skipbelow=4pt
]{ansbox}

\newmdenv[
  backgroundcolor=white,
  linecolor=darkgray,
  linewidth=0.8pt,
  leftline=true, rightline=false, topline=false, bottomline=false,
  leftmargin=0pt, rightmargin=0pt,
  innerleftmargin=8pt, innerrightmargin=6pt,
  innertopmargin=4pt, innerbottommargin=4pt,
  skipabove=4pt, skipbelow=6pt
]{invbox}

\newcommand{\ans}[1]{%
  \begin{ansbox}
    \textbf{Answer:} #1
  \end{ansbox}}

\newcommand{\method}[1]{%
  {\small\textbf{Method:} #1}\par}

\newcommand{\inv}[1]{%
  \begin{invbox}
    \textbf{\small Investigate further:} {\small #1}
  \end{invbox}}

% ── Misc helpers ────────────────────────────────────────────────────────────────
\newcommand{\qn}[1]{\textbf{#1.}\;}

% Mistake-linking: attach a Top 5 Patterns / Common Traps bullet to the question
% label(s) in this sheet where it actually shows up, e.g. \seealso{B6, D2}
\newcommand{\seealso}[1]{\hfill\textit{\footnotesize(see #1)}}

% ── Graph options macro ─────────────────────────────────────────────────────────
% Usage: \GraphOptions{tikz code for A}{tikz code for B}{tikz code for C}{tikz code for D}
\newcommand{\GraphOptions}[4]{%
  \begin{center}
    \begin{tabular}{cc}
      \textbf{A)} & \textbf{B)} \\
      #1 & #2 \\[10pt]
      \textbf{C)} & \textbf{D)} \\
      #3 & #4
    \end{tabular}
  \end{center}
}

% ── Closing motivational rule + cross-promo footer (sheets only) ────────────────
% Usage: \SpeedClosing{"Quote text."}
\newcommand{\SpeedClosing}[1]{%
  \vspace{20pt}
  \begin{center}
  \rule{0.6\textwidth}{0.4pt}\\[4pt]
  {\small\itshape #1}\\[4pt]
  \rule{0.6\textwidth}{0.4pt}
  \end{center}
  \begin{center}
  {\small Found a mistake or want to contribute a sheet?}\\
  {\small Visit the open-source project at \href{https://github.com/The-CerealDev/speed-maths}{github.com/The-CerealDev/speed-maths}}
  \end{center}
}
 and before \begin{document}):
%   \SpeedHeader{Algebra}{3}
% First arg  = pillar name shown as "Speed <Pillar> --- Daily Drill"
% Second arg = worksheet number
\pagestyle{fancy}
\newcommand{\SpeedHeader}[2]{%
  \fancyhf{}%
  \renewcommand{\headrulewidth}{0.4pt}%
  \setlength{\headheight}{14pt}%
  \fancyhead[L]{\small\textbf{Speed #1 --- Daily Drill}}%
  \fancyhead[R]{\small Worksheet \##2}%
  \fancyfoot[L]{\small \SpeedExamLine}%
  \fancyfoot[C]{\small\thepage}%
  \fancyfoot[R]{\small \SpeedCredit}%
  \renewcommand{\footrulewidth}{0.4pt}%
}

% ── Title block (used at the top of every sheet AND answer file) ───────────────
% Usage: \SpeedTitleBlock{Daily Algebra Drill \#3}{Jane Doe}
% %% AGENTS: When generating a new sheet, ask the human contributor for the name they want credited in argument 2.
% For answer files, pass the "--- Answers \& Investigations" suffix in the arg.
\newcommand{\SpeedTitleBlock}[2]{%
  \begin{center}
    {\LARGE\bfseries #1}\\[4pt]
    {\normalsize\itshape Sheet by #2}\\[6pt]
    \hrule\vspace{4pt}
    \begin{tabular}{llcc}
      \toprule
      \textbf{Section} & \textbf{Topic} & \textbf{Questions} & \textbf{Target time}\\
      \midrule
      A & Rapid Recognition          & 10 & 2 min 30 sec \\
      B & Manipulation Drills        & 10 & 8 min        \\
      C & Substitution \& Structure  & 8  & 10 min       \\
      D & Challenge                  & 5  & 15 min       \\
      \bottomrule
    \end{tabular}
    \vspace{4pt}
    \hrule
  \end{center}
}

% ── Sheet metadata: topic + the day's new tools ────────────────────────────────
% Usage (immediately after \SpeedTitleBlock, sheets only):
%   \SpeedMeta{Expanding, factorising, and the identity toolkit}
%             {difference of squares, sum/product form, completing the square}
%
% Arg 1 = the sheet's topic in one line. The website lists this instead of
%         "Sheet 03", so write the thing a reader would actually search for.
% Arg 2 = comma-separated, at most five: the tools this sheet introduces that
%         earlier sheets in the pillar did not. These become the website's tags.
%         Leave out anything the reader already met — the tags are meant to say
%         what is new today, not to summarise the sheet.
%
% Both are parsed straight out of this file by tools/build_website.py, so the
% sheet stays the single source of truth and the site cannot drift from it.
%
% This renders NOTHING. It is a declaration for the build script, not a piece of
% the sheet's layout: adding it to a sheet must not change that sheet's PDF, so
% the 25 existing sheets can be annotated without a single one of them being
% reissued. If the toolkit should also appear on the page, that is a separate
% decision about the sheet design — make it deliberately, here, not by leaving
% this macro printing something nobody asked it to.
\newcommand{\SpeedMeta}[2]{}

% ── Answer / Method / Investigate-further boxes (answer files only) ────────────
\newmdenv[
  backgroundcolor=lightgray,
  linecolor=medgray,
  linewidth=0.5pt,
  leftmargin=0pt, rightmargin=0pt,
  innerleftmargin=8pt, innerrightmargin=8pt,
  innertopmargin=5pt, innerbottommargin=5pt,
  skipabove=4pt, skipbelow=4pt
]{ansbox}

\newmdenv[
  backgroundcolor=white,
  linecolor=darkgray,
  linewidth=0.8pt,
  leftline=true, rightline=false, topline=false, bottomline=false,
  leftmargin=0pt, rightmargin=0pt,
  innerleftmargin=8pt, innerrightmargin=6pt,
  innertopmargin=4pt, innerbottommargin=4pt,
  skipabove=4pt, skipbelow=6pt
]{invbox}

\newcommand{\ans}[1]{%
  \begin{ansbox}
    \textbf{Answer:} #1
  \end{ansbox}}

\newcommand{\method}[1]{%
  {\small\textbf{Method:} #1}\par}

\newcommand{\inv}[1]{%
  \begin{invbox}
    \textbf{\small Investigate further:} {\small #1}
  \end{invbox}}

% ── Misc helpers ────────────────────────────────────────────────────────────────
\newcommand{\qn}[1]{\textbf{#1.}\;}

% Mistake-linking: attach a Top 5 Patterns / Common Traps bullet to the question
% label(s) in this sheet where it actually shows up, e.g. \seealso{B6, D2}
\newcommand{\seealso}[1]{\hfill\textit{\footnotesize(see #1)}}

% ── Graph options macro ─────────────────────────────────────────────────────────
% Usage: \GraphOptions{tikz code for A}{tikz code for B}{tikz code for C}{tikz code for D}
\newcommand{\GraphOptions}[4]{%
  \begin{center}
    \begin{tabular}{cc}
      \textbf{A)} & \textbf{B)} \\
      #1 & #2 \\[10pt]
      \textbf{C)} & \textbf{D)} \\
      #3 & #4
    \end{tabular}
  \end{center}
}

% ── Closing motivational rule + cross-promo footer (sheets only) ────────────────
% Usage: \SpeedClosing{"Quote text."}
\newcommand{\SpeedClosing}[1]{%
  \vspace{20pt}
  \begin{center}
  \rule{0.6\textwidth}{0.4pt}\\[4pt]
  {\small\itshape #1}\\[4pt]
  \rule{0.6\textwidth}{0.4pt}
  \end{center}
  \begin{center}
  {\small Found a mistake or want to contribute a sheet?}\\
  {\small Visit the open-source project at \href{https://github.com/The-CerealDev/speed-maths}{github.com/The-CerealDev/speed-maths}}
  \end{center}
}

% Editing THIS file changes the look of every sheet/answer across every pillar.
% ══════════════════════════════════════════════════════════════════════════════

\usepackage[a4paper, top=25mm, bottom=25mm, left=22mm, right=22mm]{geometry}
\usepackage{amsmath, amssymb}
\usepackage{enumitem}
\usepackage{booktabs}
\usepackage{array}
\usepackage{parskip}
\usepackage{microtype}
\usepackage{fancyhdr}
\usepackage{titlesec}
\usepackage{mdframed}
\usepackage{xcolor}
\usepackage[hidelinks]{hyperref}
\usepackage{graphicx}
\usepackage{tikz}
\usepackage{pgfplots}
\pgfplotsset{compat=1.18}

% ── Colours ───────────────────────────────────────────────────────────────────
\definecolor{darkgray}{gray}{0.15}
\definecolor{medgray}{gray}{0.45}
\definecolor{lightgray}{gray}{0.93}
\definecolor{boxgray}{gray}{0.88}

% ── Section formatting ──────────────────────────────────────────────────────────
\titleformat{\section}[block]
  {\normalfont\large\bfseries}{}
  {0pt}{}[\vspace{2pt}\hrule\vspace{6pt}]
\titlespacing*{\section}{0pt}{14pt}{4pt}

% ── List formatting ─────────────────────────────────────────────────────────────
\setlist[enumerate,1]{label=\textbf{\arabic*.}, leftmargin=2em, itemsep=10pt}

% ── Branding constants (edit here, not per-file) ────────────────────────────────
\newcommand{\SpeedExamLine}{\textbf{Speed Maths:} Competition \& Exam Prep}
\newcommand{\SpeedCredit}{Series by \href{https://www.linkedin.com/in/cerealdev/}{David Akinyele-Aje}}

% ── Header / footer ──────────────────────────────────────────────────────────────
% Usage (once, after % main (standalone preamble for editing)
% vim: nowrap: spell:
% ══════════════════════════════════════════════════════════════════════════════
% Speed Maths --- shared preamble
% Included by every sheet and answer file via:  % main (standalone preamble for editing)
% vim: nowrap: spell:
% ══════════════════════════════════════════════════════════════════════════════
% Speed Maths --- shared preamble
% Included by every sheet and answer file via:  \input{../../shared/preamble}
% Editing THIS file changes the look of every sheet/answer across every pillar.
% ══════════════════════════════════════════════════════════════════════════════

\usepackage[a4paper, top=25mm, bottom=25mm, left=22mm, right=22mm]{geometry}
\usepackage{amsmath, amssymb}
\usepackage{enumitem}
\usepackage{booktabs}
\usepackage{array}
\usepackage{parskip}
\usepackage{microtype}
\usepackage{fancyhdr}
\usepackage{titlesec}
\usepackage{mdframed}
\usepackage{xcolor}
\usepackage[hidelinks]{hyperref}
\usepackage{graphicx}
\usepackage{tikz}
\usepackage{pgfplots}
\pgfplotsset{compat=1.18}

% ── Colours ───────────────────────────────────────────────────────────────────
\definecolor{darkgray}{gray}{0.15}
\definecolor{medgray}{gray}{0.45}
\definecolor{lightgray}{gray}{0.93}
\definecolor{boxgray}{gray}{0.88}

% ── Section formatting ──────────────────────────────────────────────────────────
\titleformat{\section}[block]
  {\normalfont\large\bfseries}{}
  {0pt}{}[\vspace{2pt}\hrule\vspace{6pt}]
\titlespacing*{\section}{0pt}{14pt}{4pt}

% ── List formatting ─────────────────────────────────────────────────────────────
\setlist[enumerate,1]{label=\textbf{\arabic*.}, leftmargin=2em, itemsep=10pt}

% ── Branding constants (edit here, not per-file) ────────────────────────────────
\newcommand{\SpeedExamLine}{\textbf{Speed Maths:} Competition \& Exam Prep}
\newcommand{\SpeedCredit}{Series by \href{https://www.linkedin.com/in/cerealdev/}{David Akinyele-Aje}}

% ── Header / footer ──────────────────────────────────────────────────────────────
% Usage (once, after \input{../../shared/preamble} and before \begin{document}):
%   \SpeedHeader{Algebra}{3}
% First arg  = pillar name shown as "Speed <Pillar> --- Daily Drill"
% Second arg = worksheet number
\pagestyle{fancy}
\newcommand{\SpeedHeader}[2]{%
  \fancyhf{}%
  \renewcommand{\headrulewidth}{0.4pt}%
  \setlength{\headheight}{14pt}%
  \fancyhead[L]{\small\textbf{Speed #1 --- Daily Drill}}%
  \fancyhead[R]{\small Worksheet \##2}%
  \fancyfoot[L]{\small \SpeedExamLine}%
  \fancyfoot[C]{\small\thepage}%
  \fancyfoot[R]{\small \SpeedCredit}%
  \renewcommand{\footrulewidth}{0.4pt}%
}

% ── Title block (used at the top of every sheet AND answer file) ───────────────
% Usage: \SpeedTitleBlock{Daily Algebra Drill \#3}{Jane Doe}
% %% AGENTS: When generating a new sheet, ask the human contributor for the name they want credited in argument 2.
% For answer files, pass the "--- Answers \& Investigations" suffix in the arg.
\newcommand{\SpeedTitleBlock}[2]{%
  \begin{center}
    {\LARGE\bfseries #1}\\[4pt]
    {\normalsize\itshape Sheet by #2}\\[6pt]
    \hrule\vspace{4pt}
    \begin{tabular}{llcc}
      \toprule
      \textbf{Section} & \textbf{Topic} & \textbf{Questions} & \textbf{Target time}\\
      \midrule
      A & Rapid Recognition          & 10 & 2 min 30 sec \\
      B & Manipulation Drills        & 10 & 8 min        \\
      C & Substitution \& Structure  & 8  & 10 min       \\
      D & Challenge                  & 5  & 15 min       \\
      \bottomrule
    \end{tabular}
    \vspace{4pt}
    \hrule
  \end{center}
}

% ── Sheet metadata: topic + the day's new tools ────────────────────────────────
% Usage (immediately after \SpeedTitleBlock, sheets only):
%   \SpeedMeta{Expanding, factorising, and the identity toolkit}
%             {difference of squares, sum/product form, completing the square}
%
% Arg 1 = the sheet's topic in one line. The website lists this instead of
%         "Sheet 03", so write the thing a reader would actually search for.
% Arg 2 = comma-separated, at most five: the tools this sheet introduces that
%         earlier sheets in the pillar did not. These become the website's tags.
%         Leave out anything the reader already met — the tags are meant to say
%         what is new today, not to summarise the sheet.
%
% Both are parsed straight out of this file by tools/build_website.py, so the
% sheet stays the single source of truth and the site cannot drift from it.
%
% This renders NOTHING. It is a declaration for the build script, not a piece of
% the sheet's layout: adding it to a sheet must not change that sheet's PDF, so
% the 25 existing sheets can be annotated without a single one of them being
% reissued. If the toolkit should also appear on the page, that is a separate
% decision about the sheet design — make it deliberately, here, not by leaving
% this macro printing something nobody asked it to.
\newcommand{\SpeedMeta}[2]{}

% ── Answer / Method / Investigate-further boxes (answer files only) ────────────
\newmdenv[
  backgroundcolor=lightgray,
  linecolor=medgray,
  linewidth=0.5pt,
  leftmargin=0pt, rightmargin=0pt,
  innerleftmargin=8pt, innerrightmargin=8pt,
  innertopmargin=5pt, innerbottommargin=5pt,
  skipabove=4pt, skipbelow=4pt
]{ansbox}

\newmdenv[
  backgroundcolor=white,
  linecolor=darkgray,
  linewidth=0.8pt,
  leftline=true, rightline=false, topline=false, bottomline=false,
  leftmargin=0pt, rightmargin=0pt,
  innerleftmargin=8pt, innerrightmargin=6pt,
  innertopmargin=4pt, innerbottommargin=4pt,
  skipabove=4pt, skipbelow=6pt
]{invbox}

\newcommand{\ans}[1]{%
  \begin{ansbox}
    \textbf{Answer:} #1
  \end{ansbox}}

\newcommand{\method}[1]{%
  {\small\textbf{Method:} #1}\par}

\newcommand{\inv}[1]{%
  \begin{invbox}
    \textbf{\small Investigate further:} {\small #1}
  \end{invbox}}

% ── Misc helpers ────────────────────────────────────────────────────────────────
\newcommand{\qn}[1]{\textbf{#1.}\;}

% Mistake-linking: attach a Top 5 Patterns / Common Traps bullet to the question
% label(s) in this sheet where it actually shows up, e.g. \seealso{B6, D2}
\newcommand{\seealso}[1]{\hfill\textit{\footnotesize(see #1)}}

% ── Graph options macro ─────────────────────────────────────────────────────────
% Usage: \GraphOptions{tikz code for A}{tikz code for B}{tikz code for C}{tikz code for D}
\newcommand{\GraphOptions}[4]{%
  \begin{center}
    \begin{tabular}{cc}
      \textbf{A)} & \textbf{B)} \\
      #1 & #2 \\[10pt]
      \textbf{C)} & \textbf{D)} \\
      #3 & #4
    \end{tabular}
  \end{center}
}

% ── Closing motivational rule + cross-promo footer (sheets only) ────────────────
% Usage: \SpeedClosing{"Quote text."}
\newcommand{\SpeedClosing}[1]{%
  \vspace{20pt}
  \begin{center}
  \rule{0.6\textwidth}{0.4pt}\\[4pt]
  {\small\itshape #1}\\[4pt]
  \rule{0.6\textwidth}{0.4pt}
  \end{center}
  \begin{center}
  {\small Found a mistake or want to contribute a sheet?}\\
  {\small Visit the open-source project at \href{https://github.com/The-CerealDev/speed-maths}{github.com/The-CerealDev/speed-maths}}
  \end{center}
}

% Editing THIS file changes the look of every sheet/answer across every pillar.
% ══════════════════════════════════════════════════════════════════════════════

\usepackage[a4paper, top=25mm, bottom=25mm, left=22mm, right=22mm]{geometry}
\usepackage{amsmath, amssymb}
\usepackage{enumitem}
\usepackage{booktabs}
\usepackage{array}
\usepackage{parskip}
\usepackage{microtype}
\usepackage{fancyhdr}
\usepackage{titlesec}
\usepackage{mdframed}
\usepackage{xcolor}
\usepackage[hidelinks]{hyperref}
\usepackage{graphicx}
\usepackage{tikz}
\usepackage{pgfplots}
\pgfplotsset{compat=1.18}

% ── Colours ───────────────────────────────────────────────────────────────────
\definecolor{darkgray}{gray}{0.15}
\definecolor{medgray}{gray}{0.45}
\definecolor{lightgray}{gray}{0.93}
\definecolor{boxgray}{gray}{0.88}

% ── Section formatting ──────────────────────────────────────────────────────────
\titleformat{\section}[block]
  {\normalfont\large\bfseries}{}
  {0pt}{}[\vspace{2pt}\hrule\vspace{6pt}]
\titlespacing*{\section}{0pt}{14pt}{4pt}

% ── List formatting ─────────────────────────────────────────────────────────────
\setlist[enumerate,1]{label=\textbf{\arabic*.}, leftmargin=2em, itemsep=10pt}

% ── Branding constants (edit here, not per-file) ────────────────────────────────
\newcommand{\SpeedExamLine}{\textbf{Speed Maths:} Competition \& Exam Prep}
\newcommand{\SpeedCredit}{Series by \href{https://www.linkedin.com/in/cerealdev/}{David Akinyele-Aje}}

% ── Header / footer ──────────────────────────────────────────────────────────────
% Usage (once, after % main (standalone preamble for editing)
% vim: nowrap: spell:
% ══════════════════════════════════════════════════════════════════════════════
% Speed Maths --- shared preamble
% Included by every sheet and answer file via:  \input{../../shared/preamble}
% Editing THIS file changes the look of every sheet/answer across every pillar.
% ══════════════════════════════════════════════════════════════════════════════

\usepackage[a4paper, top=25mm, bottom=25mm, left=22mm, right=22mm]{geometry}
\usepackage{amsmath, amssymb}
\usepackage{enumitem}
\usepackage{booktabs}
\usepackage{array}
\usepackage{parskip}
\usepackage{microtype}
\usepackage{fancyhdr}
\usepackage{titlesec}
\usepackage{mdframed}
\usepackage{xcolor}
\usepackage[hidelinks]{hyperref}
\usepackage{graphicx}
\usepackage{tikz}
\usepackage{pgfplots}
\pgfplotsset{compat=1.18}

% ── Colours ───────────────────────────────────────────────────────────────────
\definecolor{darkgray}{gray}{0.15}
\definecolor{medgray}{gray}{0.45}
\definecolor{lightgray}{gray}{0.93}
\definecolor{boxgray}{gray}{0.88}

% ── Section formatting ──────────────────────────────────────────────────────────
\titleformat{\section}[block]
  {\normalfont\large\bfseries}{}
  {0pt}{}[\vspace{2pt}\hrule\vspace{6pt}]
\titlespacing*{\section}{0pt}{14pt}{4pt}

% ── List formatting ─────────────────────────────────────────────────────────────
\setlist[enumerate,1]{label=\textbf{\arabic*.}, leftmargin=2em, itemsep=10pt}

% ── Branding constants (edit here, not per-file) ────────────────────────────────
\newcommand{\SpeedExamLine}{\textbf{Speed Maths:} Competition \& Exam Prep}
\newcommand{\SpeedCredit}{Series by \href{https://www.linkedin.com/in/cerealdev/}{David Akinyele-Aje}}

% ── Header / footer ──────────────────────────────────────────────────────────────
% Usage (once, after \input{../../shared/preamble} and before \begin{document}):
%   \SpeedHeader{Algebra}{3}
% First arg  = pillar name shown as "Speed <Pillar> --- Daily Drill"
% Second arg = worksheet number
\pagestyle{fancy}
\newcommand{\SpeedHeader}[2]{%
  \fancyhf{}%
  \renewcommand{\headrulewidth}{0.4pt}%
  \setlength{\headheight}{14pt}%
  \fancyhead[L]{\small\textbf{Speed #1 --- Daily Drill}}%
  \fancyhead[R]{\small Worksheet \##2}%
  \fancyfoot[L]{\small \SpeedExamLine}%
  \fancyfoot[C]{\small\thepage}%
  \fancyfoot[R]{\small \SpeedCredit}%
  \renewcommand{\footrulewidth}{0.4pt}%
}

% ── Title block (used at the top of every sheet AND answer file) ───────────────
% Usage: \SpeedTitleBlock{Daily Algebra Drill \#3}{Jane Doe}
% %% AGENTS: When generating a new sheet, ask the human contributor for the name they want credited in argument 2.
% For answer files, pass the "--- Answers \& Investigations" suffix in the arg.
\newcommand{\SpeedTitleBlock}[2]{%
  \begin{center}
    {\LARGE\bfseries #1}\\[4pt]
    {\normalsize\itshape Sheet by #2}\\[6pt]
    \hrule\vspace{4pt}
    \begin{tabular}{llcc}
      \toprule
      \textbf{Section} & \textbf{Topic} & \textbf{Questions} & \textbf{Target time}\\
      \midrule
      A & Rapid Recognition          & 10 & 2 min 30 sec \\
      B & Manipulation Drills        & 10 & 8 min        \\
      C & Substitution \& Structure  & 8  & 10 min       \\
      D & Challenge                  & 5  & 15 min       \\
      \bottomrule
    \end{tabular}
    \vspace{4pt}
    \hrule
  \end{center}
}

% ── Sheet metadata: topic + the day's new tools ────────────────────────────────
% Usage (immediately after \SpeedTitleBlock, sheets only):
%   \SpeedMeta{Expanding, factorising, and the identity toolkit}
%             {difference of squares, sum/product form, completing the square}
%
% Arg 1 = the sheet's topic in one line. The website lists this instead of
%         "Sheet 03", so write the thing a reader would actually search for.
% Arg 2 = comma-separated, at most five: the tools this sheet introduces that
%         earlier sheets in the pillar did not. These become the website's tags.
%         Leave out anything the reader already met — the tags are meant to say
%         what is new today, not to summarise the sheet.
%
% Both are parsed straight out of this file by tools/build_website.py, so the
% sheet stays the single source of truth and the site cannot drift from it.
%
% This renders NOTHING. It is a declaration for the build script, not a piece of
% the sheet's layout: adding it to a sheet must not change that sheet's PDF, so
% the 25 existing sheets can be annotated without a single one of them being
% reissued. If the toolkit should also appear on the page, that is a separate
% decision about the sheet design — make it deliberately, here, not by leaving
% this macro printing something nobody asked it to.
\newcommand{\SpeedMeta}[2]{}

% ── Answer / Method / Investigate-further boxes (answer files only) ────────────
\newmdenv[
  backgroundcolor=lightgray,
  linecolor=medgray,
  linewidth=0.5pt,
  leftmargin=0pt, rightmargin=0pt,
  innerleftmargin=8pt, innerrightmargin=8pt,
  innertopmargin=5pt, innerbottommargin=5pt,
  skipabove=4pt, skipbelow=4pt
]{ansbox}

\newmdenv[
  backgroundcolor=white,
  linecolor=darkgray,
  linewidth=0.8pt,
  leftline=true, rightline=false, topline=false, bottomline=false,
  leftmargin=0pt, rightmargin=0pt,
  innerleftmargin=8pt, innerrightmargin=6pt,
  innertopmargin=4pt, innerbottommargin=4pt,
  skipabove=4pt, skipbelow=6pt
]{invbox}

\newcommand{\ans}[1]{%
  \begin{ansbox}
    \textbf{Answer:} #1
  \end{ansbox}}

\newcommand{\method}[1]{%
  {\small\textbf{Method:} #1}\par}

\newcommand{\inv}[1]{%
  \begin{invbox}
    \textbf{\small Investigate further:} {\small #1}
  \end{invbox}}

% ── Misc helpers ────────────────────────────────────────────────────────────────
\newcommand{\qn}[1]{\textbf{#1.}\;}

% Mistake-linking: attach a Top 5 Patterns / Common Traps bullet to the question
% label(s) in this sheet where it actually shows up, e.g. \seealso{B6, D2}
\newcommand{\seealso}[1]{\hfill\textit{\footnotesize(see #1)}}

% ── Graph options macro ─────────────────────────────────────────────────────────
% Usage: \GraphOptions{tikz code for A}{tikz code for B}{tikz code for C}{tikz code for D}
\newcommand{\GraphOptions}[4]{%
  \begin{center}
    \begin{tabular}{cc}
      \textbf{A)} & \textbf{B)} \\
      #1 & #2 \\[10pt]
      \textbf{C)} & \textbf{D)} \\
      #3 & #4
    \end{tabular}
  \end{center}
}

% ── Closing motivational rule + cross-promo footer (sheets only) ────────────────
% Usage: \SpeedClosing{"Quote text."}
\newcommand{\SpeedClosing}[1]{%
  \vspace{20pt}
  \begin{center}
  \rule{0.6\textwidth}{0.4pt}\\[4pt]
  {\small\itshape #1}\\[4pt]
  \rule{0.6\textwidth}{0.4pt}
  \end{center}
  \begin{center}
  {\small Found a mistake or want to contribute a sheet?}\\
  {\small Visit the open-source project at \href{https://github.com/The-CerealDev/speed-maths}{github.com/The-CerealDev/speed-maths}}
  \end{center}
}
 and before \begin{document}):
%   \SpeedHeader{Algebra}{3}
% First arg  = pillar name shown as "Speed <Pillar> --- Daily Drill"
% Second arg = worksheet number
\pagestyle{fancy}
\newcommand{\SpeedHeader}[2]{%
  \fancyhf{}%
  \renewcommand{\headrulewidth}{0.4pt}%
  \setlength{\headheight}{14pt}%
  \fancyhead[L]{\small\textbf{Speed #1 --- Daily Drill}}%
  \fancyhead[R]{\small Worksheet \##2}%
  \fancyfoot[L]{\small \SpeedExamLine}%
  \fancyfoot[C]{\small\thepage}%
  \fancyfoot[R]{\small \SpeedCredit}%
  \renewcommand{\footrulewidth}{0.4pt}%
}

% ── Title block (used at the top of every sheet AND answer file) ───────────────
% Usage: \SpeedTitleBlock{Daily Algebra Drill \#3}{Jane Doe}
% %% AGENTS: When generating a new sheet, ask the human contributor for the name they want credited in argument 2.
% For answer files, pass the "--- Answers \& Investigations" suffix in the arg.
\newcommand{\SpeedTitleBlock}[2]{%
  \begin{center}
    {\LARGE\bfseries #1}\\[4pt]
    {\normalsize\itshape Sheet by #2}\\[6pt]
    \hrule\vspace{4pt}
    \begin{tabular}{llcc}
      \toprule
      \textbf{Section} & \textbf{Topic} & \textbf{Questions} & \textbf{Target time}\\
      \midrule
      A & Rapid Recognition          & 10 & 2 min 30 sec \\
      B & Manipulation Drills        & 10 & 8 min        \\
      C & Substitution \& Structure  & 8  & 10 min       \\
      D & Challenge                  & 5  & 15 min       \\
      \bottomrule
    \end{tabular}
    \vspace{4pt}
    \hrule
  \end{center}
}

% ── Sheet metadata: topic + the day's new tools ────────────────────────────────
% Usage (immediately after \SpeedTitleBlock, sheets only):
%   \SpeedMeta{Expanding, factorising, and the identity toolkit}
%             {difference of squares, sum/product form, completing the square}
%
% Arg 1 = the sheet's topic in one line. The website lists this instead of
%         "Sheet 03", so write the thing a reader would actually search for.
% Arg 2 = comma-separated, at most five: the tools this sheet introduces that
%         earlier sheets in the pillar did not. These become the website's tags.
%         Leave out anything the reader already met — the tags are meant to say
%         what is new today, not to summarise the sheet.
%
% Both are parsed straight out of this file by tools/build_website.py, so the
% sheet stays the single source of truth and the site cannot drift from it.
%
% This renders NOTHING. It is a declaration for the build script, not a piece of
% the sheet's layout: adding it to a sheet must not change that sheet's PDF, so
% the 25 existing sheets can be annotated without a single one of them being
% reissued. If the toolkit should also appear on the page, that is a separate
% decision about the sheet design — make it deliberately, here, not by leaving
% this macro printing something nobody asked it to.
\newcommand{\SpeedMeta}[2]{}

% ── Answer / Method / Investigate-further boxes (answer files only) ────────────
\newmdenv[
  backgroundcolor=lightgray,
  linecolor=medgray,
  linewidth=0.5pt,
  leftmargin=0pt, rightmargin=0pt,
  innerleftmargin=8pt, innerrightmargin=8pt,
  innertopmargin=5pt, innerbottommargin=5pt,
  skipabove=4pt, skipbelow=4pt
]{ansbox}

\newmdenv[
  backgroundcolor=white,
  linecolor=darkgray,
  linewidth=0.8pt,
  leftline=true, rightline=false, topline=false, bottomline=false,
  leftmargin=0pt, rightmargin=0pt,
  innerleftmargin=8pt, innerrightmargin=6pt,
  innertopmargin=4pt, innerbottommargin=4pt,
  skipabove=4pt, skipbelow=6pt
]{invbox}

\newcommand{\ans}[1]{%
  \begin{ansbox}
    \textbf{Answer:} #1
  \end{ansbox}}

\newcommand{\method}[1]{%
  {\small\textbf{Method:} #1}\par}

\newcommand{\inv}[1]{%
  \begin{invbox}
    \textbf{\small Investigate further:} {\small #1}
  \end{invbox}}

% ── Misc helpers ────────────────────────────────────────────────────────────────
\newcommand{\qn}[1]{\textbf{#1.}\;}

% Mistake-linking: attach a Top 5 Patterns / Common Traps bullet to the question
% label(s) in this sheet where it actually shows up, e.g. \seealso{B6, D2}
\newcommand{\seealso}[1]{\hfill\textit{\footnotesize(see #1)}}

% ── Graph options macro ─────────────────────────────────────────────────────────
% Usage: \GraphOptions{tikz code for A}{tikz code for B}{tikz code for C}{tikz code for D}
\newcommand{\GraphOptions}[4]{%
  \begin{center}
    \begin{tabular}{cc}
      \textbf{A)} & \textbf{B)} \\
      #1 & #2 \\[10pt]
      \textbf{C)} & \textbf{D)} \\
      #3 & #4
    \end{tabular}
  \end{center}
}

% ── Closing motivational rule + cross-promo footer (sheets only) ────────────────
% Usage: \SpeedClosing{"Quote text."}
\newcommand{\SpeedClosing}[1]{%
  \vspace{20pt}
  \begin{center}
  \rule{0.6\textwidth}{0.4pt}\\[4pt]
  {\small\itshape #1}\\[4pt]
  \rule{0.6\textwidth}{0.4pt}
  \end{center}
  \begin{center}
  {\small Found a mistake or want to contribute a sheet?}\\
  {\small Visit the open-source project at \href{https://github.com/The-CerealDev/speed-maths}{github.com/The-CerealDev/speed-maths}}
  \end{center}
}
 and before \begin{document}):
%   \SpeedHeader{Algebra}{3}
% First arg  = pillar name shown as "Speed <Pillar> --- Daily Drill"
% Second arg = worksheet number
\pagestyle{fancy}
\newcommand{\SpeedHeader}[2]{%
  \fancyhf{}%
  \renewcommand{\headrulewidth}{0.4pt}%
  \setlength{\headheight}{14pt}%
  \fancyhead[L]{\small\textbf{Speed #1 --- Daily Drill}}%
  \fancyhead[R]{\small Worksheet \##2}%
  \fancyfoot[L]{\small \SpeedExamLine}%
  \fancyfoot[C]{\small\thepage}%
  \fancyfoot[R]{\small \SpeedCredit}%
  \renewcommand{\footrulewidth}{0.4pt}%
}

% ── Title block (used at the top of every sheet AND answer file) ───────────────
% Usage: \SpeedTitleBlock{Daily Algebra Drill \#3}{Jane Doe}
% %% AGENTS: When generating a new sheet, ask the human contributor for the name they want credited in argument 2.
% For answer files, pass the "--- Answers \& Investigations" suffix in the arg.
\newcommand{\SpeedTitleBlock}[2]{%
  \begin{center}
    {\LARGE\bfseries #1}\\[4pt]
    {\normalsize\itshape Sheet by #2}\\[6pt]
    \hrule\vspace{4pt}
    \begin{tabular}{llcc}
      \toprule
      \textbf{Section} & \textbf{Topic} & \textbf{Questions} & \textbf{Target time}\\
      \midrule
      A & Rapid Recognition          & 10 & 2 min 30 sec \\
      B & Manipulation Drills        & 10 & 8 min        \\
      C & Substitution \& Structure  & 8  & 10 min       \\
      D & Challenge                  & 5  & 15 min       \\
      \bottomrule
    \end{tabular}
    \vspace{4pt}
    \hrule
  \end{center}
}

% ── Sheet metadata: topic + the day's new tools ────────────────────────────────
% Usage (immediately after \SpeedTitleBlock, sheets only):
%   \SpeedMeta{Expanding, factorising, and the identity toolkit}
%             {difference of squares, sum/product form, completing the square}
%
% Arg 1 = the sheet's topic in one line. The website lists this instead of
%         "Sheet 03", so write the thing a reader would actually search for.
% Arg 2 = comma-separated, at most five: the tools this sheet introduces that
%         earlier sheets in the pillar did not. These become the website's tags.
%         Leave out anything the reader already met — the tags are meant to say
%         what is new today, not to summarise the sheet.
%
% Both are parsed straight out of this file by tools/build_website.py, so the
% sheet stays the single source of truth and the site cannot drift from it.
%
% This renders NOTHING. It is a declaration for the build script, not a piece of
% the sheet's layout: adding it to a sheet must not change that sheet's PDF, so
% the 25 existing sheets can be annotated without a single one of them being
% reissued. If the toolkit should also appear on the page, that is a separate
% decision about the sheet design — make it deliberately, here, not by leaving
% this macro printing something nobody asked it to.
\newcommand{\SpeedMeta}[2]{}

% ── Answer / Method / Investigate-further boxes (answer files only) ────────────
\newmdenv[
  backgroundcolor=lightgray,
  linecolor=medgray,
  linewidth=0.5pt,
  leftmargin=0pt, rightmargin=0pt,
  innerleftmargin=8pt, innerrightmargin=8pt,
  innertopmargin=5pt, innerbottommargin=5pt,
  skipabove=4pt, skipbelow=4pt
]{ansbox}

\newmdenv[
  backgroundcolor=white,
  linecolor=darkgray,
  linewidth=0.8pt,
  leftline=true, rightline=false, topline=false, bottomline=false,
  leftmargin=0pt, rightmargin=0pt,
  innerleftmargin=8pt, innerrightmargin=6pt,
  innertopmargin=4pt, innerbottommargin=4pt,
  skipabove=4pt, skipbelow=6pt
]{invbox}

\newcommand{\ans}[1]{%
  \begin{ansbox}
    \textbf{Answer:} #1
  \end{ansbox}}

\newcommand{\method}[1]{%
  {\small\textbf{Method:} #1}\par}

\newcommand{\inv}[1]{%
  \begin{invbox}
    \textbf{\small Investigate further:} {\small #1}
  \end{invbox}}

% ── Misc helpers ────────────────────────────────────────────────────────────────
\newcommand{\qn}[1]{\textbf{#1.}\;}

% Mistake-linking: attach a Top 5 Patterns / Common Traps bullet to the question
% label(s) in this sheet where it actually shows up, e.g. \seealso{B6, D2}
\newcommand{\seealso}[1]{\hfill\textit{\footnotesize(see #1)}}

% ── Graph options macro ─────────────────────────────────────────────────────────
% Usage: \GraphOptions{tikz code for A}{tikz code for B}{tikz code for C}{tikz code for D}
\newcommand{\GraphOptions}[4]{%
  \begin{center}
    \begin{tabular}{cc}
      \textbf{A)} & \textbf{B)} \\
      #1 & #2 \\[10pt]
      \textbf{C)} & \textbf{D)} \\
      #3 & #4
    \end{tabular}
  \end{center}
}

% ── Closing motivational rule + cross-promo footer (sheets only) ────────────────
% Usage: \SpeedClosing{"Quote text."}
\newcommand{\SpeedClosing}[1]{%
  \vspace{20pt}
  \begin{center}
  \rule{0.6\textwidth}{0.4pt}\\[4pt]
  {\small\itshape #1}\\[4pt]
  \rule{0.6\textwidth}{0.4pt}
  \end{center}
  \begin{center}
  {\small Found a mistake or want to contribute a sheet?}\\
  {\small Visit the open-source project at \href{https://github.com/The-CerealDev/speed-maths}{github.com/The-CerealDev/speed-maths}}
  \end{center}
}

\SpeedHeader{Logic}{2}

\begin{document}

\SpeedTitleBlock{Daily Logic Drill \#2 --- Answers \& Investigations}{David Akinyele-Aje}
\noindent\textit{New toolkit today: Implication, Converse, Inverse \& Contrapositive $\cdot$ Necessary vs Sufficient.}


\section*{Section A \quad Rapid Recognition}

\begin{enumerate}

%── A1 ──────────────────────────────────────────────────────────────────────────
\item For real numbers $x, y$, classify ``$x > y$'' as a condition for ``$x^3 > y^3$''.

\ans{Necessary and sufficient.}
\method{The function $f(t)=t^3$ has derivative $f'(t)=3t^2\geq0$, with $f'(t)=0$ only at $t=0$, so $f$ is strictly increasing everywhere on $\mathbb{R}$. Hence $x > y \iff x^3 > y^3$. Both directions hold.}
\inv{Unlike squaring ($x > y \not\iff x^2 > y^2$), odd powers preserve order globally across all real numbers.}

%── A2 ──────────────────────────────────────────────────────────────────────────
\item For non-zero real numbers $x, y$, classify ``$x > y$'' as a condition for ``$\dfrac{1}{x} < \dfrac{1}{y}$''.

\ans{Neither necessary nor sufficient.}
\method{Not sufficient: $x=1, y=-2 \implies x > y$, but $\frac{1}{x} = 1 > -\frac{1}{2} = \frac{1}{y}$. Not necessary: $x=-2, y=-1 \implies \frac{1}{x} = -0.5 > -1 = \frac{1}{y}$, while for $x=-1, y=2$ we have $\frac{1}{x} = -1 < \frac{1}{2} = \frac{1}{y}$ but $x < y$.}
\inv{The inequality $\frac{1}{x} < \frac{1}{y} \iff \frac{y-x}{xy} < 0$ depends entirely on the sign of the product $xy$. It only reverses when $x$ and $y$ share the same sign.}

%── A3 ──────────────────────────────────────────────────────────────────────────
\item For the quadratic equation $x^2 - 2kx + 1 = 0$, classify ``$k > 1$'' as a condition for ``the equation has two distinct positive real roots''.

\ans{Necessary and sufficient.}
\method{For two distinct real roots, the discriminant $\Delta = 4k^2 - 4 > 0 \iff k^2 > 1 \iff |k| > 1$. By Vieta's formulas, the product of roots is $1 > 0$, so the roots share the same sign. The sum of roots is $2k$. For the roots to be positive, we need sum $2k > 0 \implies k > 0$. Combining $|k| > 1$ and $k > 0$ gives $k > 1$.}
\inv{This three-part test (discriminant, product sign, sum sign) is the universal TMUA protocol for quadratic root separation.}

%── A4 ──────────────────────────────────────────────────────────────────────────
\item For real numbers $x, y$, classify ``$|x + y| = |x| + |y|$'' as a condition for ``$xy \ge 0$''.

\ans{Necessary and sufficient.}
\method{Both sides are non-negative, so square both sides: $(x+y)^2 = (|x|+|y|)^2 \iff x^2+2xy+y^2 = x^2+2|xy|+y^2 \iff xy = |xy|$. Since $xy = |xy| \iff xy \geq 0$, the condition is necessary and sufficient.}
\inv{The Triangle Inequality $|x+y| \le |x|+|y|$ achieves equality if and only if $x$ and $y$ have the same sign (or at least one is zero).}

%── A5 ──────────────────────────────────────────────────────────────────────────
\item Let $f:\mathbb{R}\to\mathbb{R}$ be differentiable. Classify ``$f'(x) > 0$ for all $x \in \mathbb{R}$'' as a condition for ``$f$ is strictly increasing on $\mathbb{R}$''.

\ans{Sufficient but not necessary.}
\method{By the Mean Value Theorem, $f'(x) > 0$ everywhere implies $f(b)-f(a) = f'(c)(b-a) > 0$ for all $b > a$, so $f$ is strictly increasing (sufficient). It is not necessary: $f(x) = x^3$ is strictly increasing on $\mathbb{R}$, yet $f'(0) = 0 \not> 0$.}
\inv{A classic TMUA distinction: strict positivity of the derivative is sufficient, but isolated stationary points do not break strict monotonicity.}

%── A6 ──────────────────────────────────────────────────────────────────────────
\item For real constants $a, b, c, d$, classify ``$ad - bc \neq 0$'' as a condition for ``the system $ax + by = 0,\ cx + dy = 0$ has only the trivial solution $(x,y)=(0,0)$''.

\ans{Necessary and sufficient.}
\method{A homogeneous $2\times2$ linear system has non-trivial solutions if and only if the lines are collinear, which occurs iff $\det = ad - bc = 0$. Hence having only the trivial solution is equivalent to $ad - bc \neq 0$.}
\inv{Determinant non-zero $\iff$ unique solution $\iff$ kernel is trivial: the foundation of linear algebra condition strength.}

%── A7 ──────────────────────────────────────────────────────────────────────────
\item For a point $(x, y) \in \mathbb{R}^2$, classify ``$x^2 + y^2 \le 1$'' as a condition for ``$|x| + |y| \le \sqrt{2}$''.

\ans{Sufficient but not necessary.}
\method{By Cauchy-Schwarz, $(|x|+|y|)^2 \le (1^2+1^2)(x^2+y^2) = 2(x^2+y^2)$. If $x^2+y^2 \le 1$, then $(|x|+|y|)^2 \le 2 \implies |x|+|y| \le \sqrt{2}$ (sufficient). Not necessary: the point $(1, 0.4)$ has $|x|+|y| = 1.4 \le \sqrt{2} \approx 1.414$, but $x^2+y^2 = 1 + 0.16 = 1.16 > 1$.}
\inv{Geometrically, the unit circle is inscribed inside the diamond $|x|+|y| \le \sqrt{2}$ with four tangent points $(\pm\frac{1}{\sqrt2}, \pm\frac{1}{\sqrt2})$. The diamond strictly contains the disk.}

%── A8 ──────────────────────────────────────────────────────────────────────────
\item For a positive integer $n$, classify ``$12 \mid n$'' as a condition for ``$4 \mid n$ and $6 \mid n$''.

\ans{Necessary and sufficient.}
\method{An integer $n$ is divisible by both $4$ and $6$ if and only if $n$ is divisible by $\operatorname{lcm}(4, 6) = 12$. Since $12 \mid n \iff (4 \mid n \land 6 \mid n)$, the condition is necessary and sufficient.}
\inv{Unlike $4 \times 6 = 24$, which is sufficient but not necessary, the LCM gives the exact biconditional.}

%── A9 ──────────────────────────────────────────────────────────────────────────
\item For a triangle with side lengths $a, b, c$, classify ``$a^2 + b^2 = c^2$'' as a condition for ``the area of the triangle is $\frac{1}{2}ab$''.

\ans{Necessary and sufficient.}
\method{The area of any triangle with sides $a, b$ and included angle $C$ is $\text{Area} = \frac{1}{2}ab\sin C$. Thus $\text{Area} = \frac{1}{2}ab \iff \sin C = 1 \iff C = 90^\circ$. By the Law of Cosines, $c^2 = a^2+b^2 - 2ab\cos C$, so $C = 90^\circ \iff a^2+b^2 = c^2$. Both directions hold.}
\inv{Pythagoras and the sine area formula link right-angles to area maximization: $\sin C \le 1$ with equality iff right-angled.}

%── A10 ─────────────────────────────────────────────────────────────────────────
\item For an integer $n$, classify ``$n$ is odd'' as a condition for ``$n^2 \equiv 1 \pmod 8$''.

\ans{Necessary and sufficient.}
\method{If $n$ is odd, $n = 2k+1$, so $n^2 = 4k(k+1)+1$. Since $k(k+1)$ is the product of consecutive integers, it is even, so $4k(k+1)$ is a multiple of $8$, giving $n^2 \equiv 1 \pmod 8$ (sufficient). Conversely, if $n$ is even, $n^2$ is even, so $n^2 \not\equiv 1 \pmod 8$ (necessary).}
\inv{Every odd square is $1 \pmod 8$ --- one of the most powerful modular invariants in TMUA Paper 2.}

\end{enumerate}

\section*{Section B \quad Manipulation Drills}

\begin{enumerate}

%── B1 ──────────────────────────────────────────────────────────────────────────
\item Determine the necessary and sufficient condition on real $k$ such that the quadratic equation $x^2 - kx + 9 = 0$ has two distinct real roots that are both strictly greater than $1$. \textit{\small(after TMUA 2021 Paper 2 Q9)}

\ans{$6 < k < 10$.}
\method{Let $f(x) = x^2 - kx + 9$. For $f$ to have two distinct real roots $> 1$: (1) $\Delta = k^2 - 36 > 0 \implies |k| > 6$. (2) The vertex of the parabola $x = k/2 > 1 \implies k > 2$. Combining with $|k| > 6$ gives $k > 6$. (3) Both roots $> 1$ requires $f(1) > 0$, so $f(1) = 1 - k + 9 = 10 - k > 0 \implies k < 10$. Together, $6 < k < 10$.}
\inv{If $k \ge 10$, one root is $\le 1$; if $k \le 6$, the roots are complex or coincident.}

%── B2 ──────────────────────────────────────────────────────────────────────────
\item Let $f:\mathbb{R}\to\mathbb{R}$ be differentiable. Consider the conditions:
\begin{enumerate}
\item[I.] $f(-x) = f(x)$ for all $x \in \mathbb{R}$.
\item[II.] $f'(-x) = -f'(x)$ for all $x \in \mathbb{R}$.
\end{enumerate}
Classify condition I as a condition for condition II. \textit{\small(after TMUA 2018 Paper 2 Q10)}

\ans{Necessary and sufficient.}
\method{If I holds, differentiating $f(-x) = f(x)$ via chain rule gives $-f'(-x) = f'(x) \implies f'(-x) = -f'(x)$, so II holds (sufficient). Conversely, if II holds, let $h(x) = f(x) - f(-x)$. Then $h'(x) = f'(x) - (-f'(-x)) = f'(x) + f'(-x) = 0$ for all $x$, so $h(x) = C$ is constant. At $x = 0$, $h(0) = f(0) - f(0) = 0 \implies C = 0$, so $f(x) = f(-x)$ for all $x$ (necessary).}
\inv{Because $h(0) = 0$ is guaranteed, the constant of integration vanishes, making evenness and odd derivative strictly equivalent on $\mathbb{R}$.}

%── B3 ──────────────────────────────────────────────────────────────────────────
\item For positive real numbers $x, y$, classify ``$\ln(x+y) < \ln x + \ln y$'' as a condition for ``$x > 1$ and $y > 1$''.

\ans{Neither necessary nor sufficient.}
\method{$\ln(x+y) < \ln(xy) \iff x+y < xy \iff xy - x - y > 0 \iff (x-1)(y-1) > 1$. Not sufficient: $x = 1.5, y = 1.5 \implies x > 1$ and $y > 1$, but $(0.5)(0.5) = 0.25 \not> 1$. Not necessary: $x = 10, y = 1.2 \implies (9)(0.2) = 1.8 > 1$, but if $x = 3, y = 0.5$ then $(2)(-0.5) = -1 < 1$. Furthermore, $x > 1$ and $y > 1$ is not necessary because both $x-1$ and $y-1$ could never be negative and multiply to $>1$ while being positive reals (since $x, y > 0 \implies x-1 > -1, y-1 > -1$). So it is neither.}
\inv{Factoring $xy - x - y = (x-1)(y-1)-1$ converts the logarithmic inequality into a hyperbola region.}

%── B4 ──────────────────────────────────────────────────────────────────────────
\item State the contrapositive of: ``If $a^2 + b^2$ is divisible by $3$, then both $a$ and $b$ are divisible by $3$'' (for integers $a, b$). Is this contrapositive true?

\ans{Contrapositive: ``If at least one of $a$ or $b$ is not divisible by $3$, then $a^2 + b^2$ is not divisible by $3$.'' True.}
\method{The contrapositive of $P \implies (Q \land R)$ is $\neg(Q \land R) \implies \neg P \equiv (\neg Q \lor \neg R) \implies \neg P$. In words: if $a$ or $b$ is not divisible by $3$, then $a^2+b^2$ is not divisible by $3$. For any integer $n$, $n^2 \equiv 0$ or $1 \pmod 3$. Thus $a^2+b^2 \equiv 0 \pmod 3$ iff $a^2 \equiv 0$ and $b^2 \equiv 0 \pmod 3 \iff 3\mid a$ and $3\mid b$. The original statement is true, so its contrapositive is true.}
\inv{De Morgan on the conclusion of an implication flips ``both'' to ``at least one fails''.}

%── B5 ──────────────────────────────────────────────────────────────────────────
\item Let $T_1$ and $T_2$ be two triangles. Classify ``$T_1$ and $T_2$ have equal perimeter and equal area'' as a condition for ``$T_1$ and $T_2$ are congruent''.

\ans{Necessary but not sufficient.}
\method{Congruent triangles have identical sides, so equal perimeter and equal area is necessary. It is not sufficient: a triangle's area and perimeter fix the semiperimeter $s$ and inradius $r = A/s$, but there exist non-congruent triangles with the same $s$ and $r$ (e.g.\ non-isometric Heron triangles or right triangles vs oblique triangles sharing $s$ and $A$).}
\inv{Two invariants (perimeter and area) are not enough to fix the three degrees of freedom of a triangle up to congruence.}

%── B6 ──────────────────────────────────────────────────────────────────────────
\item For an integer $n$, classify ``$n$ is odd'' as a condition for ``$n^3 - n$ is a multiple of $24$''.

\ans{Sufficient but not necessary.}
\method{Factor $n^3-n = (n-1)n(n+1)$. If $n$ is odd, $n-1$ and $n+1$ are consecutive even integers, so one is divisible by $4$ and the other by $2$, giving $8 \mid n^3-n$. Among three consecutive integers, one is always divisible by $3$, so $24 \mid n^3-n$ (sufficient). It is not necessary: if $n = 8$ (even), $8^3 - 8 = 504 = 24 \times 21$, which is divisible by $24$.}
\inv{A classic pitfall: $n$ odd guarantees divisibility by $24$, but certain even integers ($n$ a multiple of $24$, or $n \equiv 0 \pmod 8$) also work.}

%── B7 ──────────────────────────────────────────────────────────────────────────
\item Suppose propositions $P, Q, R$ satisfy $P \implies Q$ and $R \implies Q$. Must either $P \implies R$ or $R \implies P$ be true? Justify with a mathematical counterexample.

\ans{No. (Counterexample: let $Q$ be $x^2 \ge 0$, $P$ be $x = 1$, and $R$ be $x = -1$).}
\method{Both $x=1$ and $x=-1$ imply $x^2 \ge 0$, so $P\implies Q$ and $R\implies Q$ are both true. But $x=1 \implies x=-1$ is false, and $x=-1 \implies x=1$ is false. Thus neither $P\implies R$ nor $R\implies P$ holds.}
\inv{Two independent sufficient conditions for the same conclusion need have no logical relation to each other.}

%── B8 ──────────────────────────────────────────────────────────────────────────
\item For real $c$, determine whether ``$c < 0$'' is necessary, sufficient, both, or neither for the quartic equation $x^4 + x^2 + c = 0$ to have at least two distinct real roots.

\ans{Necessary and sufficient.}
\method{Let $u = x^2 \ge 0$. The equation becomes $u^2 + u + c = 0$. By Vieta's formulas, the product of roots in $u$ is $c$. If $c < 0$, the roots in $u$ have opposite signs, so there is exactly one positive root $u_1 > 0$, yielding two distinct real roots $x = \pm\sqrt{u_1}$ (sufficient). If $c \ge 0$, then for all $u \ge 0$, $u^2+u+c \ge c \ge 0$; the only possible real root is $u=0$ when $c=0$, which gives $x=0$ (one root, not two distinct roots). Thus $c < 0$ is also necessary.}
\inv{Substitution $u=x^2$ requires paying close attention to the domain $u \ge 0$: only positive $u$ generate pairs of real roots.}

%── B9 ──────────────────────────────────────────────────────────────────────────
\item Consider the statement: ``If $n$ is a multiple of $9$, then $n^2$ is a multiple of $27$'' (for positive integer $n$). State its converse, and determine whether the converse is true or false.

\ans{Converse: ``If $n^2$ is a multiple of $27$, then $n$ is a multiple of $9$.'' True.}
\method{The converse states $27 \mid n^2 \implies 9 \mid n$. If $27 \mid n^2$, then $3^3 \mid n^2$. Let the power of $3$ dividing $n$ be $3^k$. Then $3^{2k}$ divides $n^2$, so $2k \ge 3 \implies k \ge 2$. Thus $3^2 = 9$ divides $n$. Hence the converse is true.}
\inv{The exponent $3$ in $27 = 3^3$ forces $k \ge \lceil 3/2 \rceil = 2$, which uniquely forces divisibility by $3^2 = 9$.}

%── B10 ─────────────────────────────────────────────────────────────────────────
\item Let $f:\mathbb{R}\to\mathbb{R}$ be twice differentiable. Classify ``$f'(a) = 0$ and $f''(a) > 0$'' as a condition for ``$f$ has a local minimum at $x = a$''.

\ans{Sufficient but not necessary.}
\method{By the Second Derivative Test, $f'(a) = 0$ and $f''(a) > 0$ guarantees a local minimum (sufficient). It is not necessary: $f(x) = x^4$ has a global minimum at $x = 0$, where $f'(0) = 0$ and $f''(0) = 0 \not> 0$.}
\inv{Higher-order flat minima ($x^4, x^6$) refute necessity of strict second derivatives.}

\end{enumerate}

\section*{Section C \quad Substitution \& Structure}

\begin{enumerate}

%── C1 ──────────────────────────────────────────────────────────────────────────
\item For real numbers $x,y$, the condition ``$x+y$ is an integer and $x-y$ is an integer'' is:

\ans{A) necessary but not sufficient for ``$x$ and $y$ are both integers''.}
\method{Necessary: if $x,y$ are both integers, then $x+y$ and $x-y$ are automatically both integers. Not sufficient: $x=y=0.5$ gives $x+y=1$ and $x-y=0$, both integers, but $x,y$ themselves are not --- ruling out options B and C. Necessity alone rules out D.}
\inv{Solving $x=\frac{(x+y)+(x-y)}{2}$, $y=\frac{(x+y)-(x-y)}{2}$ shows exactly what goes wrong: $x,y$ are only forced to be integers when $x+y$ and $x-y$ have the \emph{same parity} (so the sum/difference is even) --- a condition the question's hypothesis never guarantees.}

%── C2 ──────────────────────────────────────────────────────────────────────────
\item Determine the condition on real parameters $p, q$ that is necessary and sufficient for the cubic $x^3 - 3px + q = 0$ with $p > 0$ to have three distinct real roots.

\ans{B) $4p^3 > q^2$.}
\method{Let $f(x) = x^3 - 3px + q$. Differentiating gives $f'(x) = 3x^2 - 3p = 3(x^2 - p) = 0 \implies x = \pm\sqrt{p}$. A cubic with positive leading coefficient has three distinct real roots iff its local maximum is positive and local minimum is negative: $f(-\sqrt{p}) > 0$ and $f(\sqrt{p}) < 0$, which is equivalent to $f(-\sqrt{p})f(\sqrt{p}) < 0$. We compute: $f(\sqrt{p}) = p^{3/2} - 3p^{3/2} + q = q - 2p^{3/2}$, and $f(-\sqrt{p}) = -p^{3/2} + 3p^{3/2} + q = q + 2p^{3/2}$. Their product is $(q - 2p^{3/2})(q + 2p^{3/2}) = q^2 - 4p^3 < 0 \iff 4p^3 > q^2$. This matches option B.}
\inv{This is the classic cubic discriminant $\Delta = 4p^3 - q^2 > 0$ derived entirely from calculus extrema.}

%── C3 ──────────────────────────────────────────────────────────────────────────
\item Classify ``a triangle has all three altitudes of equal length'' as a condition for ``the triangle is equilateral''. Justify using the fact that an altitude to side $a$ has length $\frac{2\times\text{Area}}{a}$.

\ans{C) necessary and sufficient.}
\method{Let the altitudes to sides $a,b,c$ be $h_a,h_b,h_c$, and let the triangle's area be $A$. Sufficient: if the triangle is equilateral, $a=b=c$, so $h_a=\frac{2A}{a}=\frac{2A}{b}=\frac{2A}{c}=h_b=h_c$ directly. Necessary: if $h_a=h_b=h_c$, then $\frac{2A}{a}=\frac{2A}{b}=\frac{2A}{c}$, and since $A\neq0$, this forces $a=b=c$ --- the triangle is equilateral.}
\inv{Both directions rely on the same identity ($h_a=\frac{2A}{a}$) read in opposite directions --- side-length equality and altitude-length equality are simply the same fact viewed through a common factor of $2A$.}

%── C4 ──────────────────────────────────────────────────────────────────────────
\item For positive integers $a,b,n$, suppose you're told only that $n$ divides $a$, or $n$ divides $b$ (possibly both). Does this pin down whether $n$ divides the product $ab$? And working the other way: does knowing $n$ divides $ab$ pin down whether $n$ divides $a$ or $n$ divides $b$? \textit{\small(after TMUA 2019 Paper 2 Q5)}

\ans{B) ``$n$ divides $a$ or $n$ divides $b$'' is sufficient but not necessary for ``$n$ divides $ab$''.}
\method{Sufficient: if $n\mid a$ (or $n\mid b$), then $n$ divides the product $ab$ directly. Not necessary: taking $n=9$, $a=3$, $b=3$ gives $ab=9$, divisible by $9$, but neither $a$ nor $b$ individually is divisible by $9$ --- ruling out options A and C. Option D is ruled out by the sufficiency direction alone.}
\inv{This only fails because $n=9=3^2$ is not squarefree --- if $n$ were prime, divisibility of a product would force divisibility of a factor (Euclid's Lemma), making the condition necessary too. The gap here is exactly why primality matters so much in Number Theory.}

%── C5 ──────────────────────────────────────────────────────────────────────────
\item For a twice-differentiable function $f:\mathbb{R}\to\mathbb{R}$, which of the following is necessary and sufficient for the midpoint convexity condition $\frac{f(a)+f(b)}{2} \ge f\left(\frac{a+b}{2}\right)$ to hold for all $a, b \in \mathbb{R}$? \textit{\small(after TMUA 2018 Paper 2 Q18)}

\ans{B) $f''(x) \ge 0$ for all $x \in \mathbb{R}$.}
\method{For a twice-differentiable function, convexity is equivalent to non-negative second derivative: $f''(x) \ge 0$ for all $x$. Option A ($f''(x) > 0$) is sufficient but not necessary: $f(x) = x^4$ is convex everywhere but $f''(0) = 0$, and linear functions $f(x) = mx+c$ satisfy equality with $f''(x) \equiv 0$. Option C and D have no relation to convexity. Hence B is necessary and sufficient.}
\inv{Midpoint convexity plus continuity implies full Jensen convexity $\sum \lambda_i f(x_i) \ge f(\sum \lambda_i x_i)$.}

%── C6 ──────────────────────────────────────────────────────────────────────────
\item Classify ``$n^2$ is even'' as a condition for ``$n$ is even''.

\ans{C) necessary and sufficient.}
\method{Sufficient: if $n^2$ is even, then $n$ cannot be odd (an odd number squared is odd), so $n$ is even. Necessary: if $n$ is even, $n^2$ is even directly.}
\inv{Both directions holding here is exactly why ``$n$ even $\iff n^2$ even'' is such a reliable tool in proofs by contradiction.}

%── C7 ──────────────────────────────────────────────────────────────────────────
\item Let $p(x)$ be a cubic polynomial with real coefficients. Which of the following conditions is necessary and sufficient for the existence of $c \in (a, b)$ such that $p'(c) = 0$? \textit{\small(after TMUA 2021 Paper 2 Q8)}

\ans{C) $p'(x) = 0$ has real roots and at least one lies in $(a, b)$.}
\method{The statement asserts the existence of a real number $c \in (a, b)$ satisfying $p'(c) = 0$. By definition, this occurs if and only if $p'(x) = 0$ has real roots and at least one root lies in the open interval $(a, b)$. Option A ($p(a)=p(b)$) is sufficient by Rolle's Theorem, but not necessary (a cubic can have a stationary point in $(a, b)$ without $p(a)=p(b)$). Option B is sufficient by the Intermediate Value Theorem, but not necessary (both $p'(a)$ and $p'(b)$ could be positive while a local min and max both lie inside $(a, b)$). Thus C is the exact necessary and sufficient condition.}
\inv{Distinguishing between a theorem's hypothesis (sufficient, like Rolle's) and the exact semantic equivalence is a core TMUA Paper 2 theme.}

%── C8 ──────────────────────────────────────────────────────────────────────────
\item State the contrapositive of ``If $a$ and $b$ are both irrational, then $ab$ is irrational.'' Is the original statement true? Give a counterexample if false.

\ans{Contrapositive: ``If $ab$ is rational, then $a$ is rational or $b$ is rational.'' The original statement is false; counterexample $a=b=\sqrt2$.}
\method{$\sqrt2$ is irrational, but $ab=\sqrt2\times\sqrt2=2$, which is rational --- disproving the original. Since the contrapositive always shares the original's truth value, the same counterexample shows the contrapositive is false too: $ab=2$ is rational, yet neither $a=\sqrt2$ nor $b=\sqrt2$ is rational.}
\inv{One counterexample disproves both a statement and its contrapositive simultaneously --- there is never a need to separately hunt for a second one.}

\end{enumerate}

\section*{Section D \quad Challenge}

\begin{enumerate}

%── D1 ──────────────────────────────────────────────────────────────────────────
\item Let $f(x)=x^2+bx+c$ with $f(0)=c$. Let $P$ be the statement: ``if $c<0$, then $f$ has two distinct real roots.'' Which of the following must be true? \textit{\small(after TMUA 2022 Paper 2 Q5)}

\ans{F) I and III only.}
\method{I ($P$): the discriminant is $b^2-4c$; if $c<0$ then $-4c>0$, so $b^2-4c > b^2 \geq 0$, meaning the discriminant is always strictly positive --- $f$ always has two distinct real roots. $P$ is true. II (converse): ``if $f$ has two distinct real roots, then $c<0$'' --- false, since $b=3,c=1$ gives discriminant $9-4=5>0$ (two distinct real roots) with $c=1\not<0$. III (contrapositive): shares $P$'s truth value, so it is true.}
\inv{A true original implication with a false converse and a true contrapositive is the quintessential TMUA implication pattern.}

%── D2 ──────────────────────────────────────────────────────────────────────────
\item Consider a finite list of $n$ real numbers (repeats allowed), sorted into non-decreasing order. Let $P$ be the statement ``$n$ is odd'', and let $Q$ be the statement ``the middle value of the sorted list equals one of the list's own entries''. $P$ is: \textit{\small(after TMUA 2022 Paper 2 Q6)}

\ans{B) sufficient but not necessary for $Q$.}
\method{Sufficient: if $n$ is odd, the list has a unique middle term when sorted, and that middle term \emph{is} the median --- a list member. Not necessary: the list $4,4$ has $n=2$ (even), and the median is $\frac{4+4}{2}=4$, which is a list member --- so $Q$ holds even though $P$ fails, ruling out options A and C. Sufficiency alone rules out D.}
\inv{Any even-length list whose two middle entries are equal gives the same kind of counterexample --- it is $n$ being odd that \emph{guarantees} $Q$, not the only way to achieve it.}

%── D3 ──────────────────────────────────────────────────────────────────────────
\item For a positive integer $n$, let $P$ be ``$n$ leaves remainder $1$ when divided by $4$'' and $Q$ be ``$n^2$ leaves remainder $1$ when divided by $8$.'' $P$ is:

\ans{B) sufficient but not necessary for $Q$.}
\method{Sufficient: if $n=4k+1$, then $n^2=16k^2+8k+1$, and modulo $8$ this is $1$ (both $16k^2$ and $8k$ are multiples of $8$). Not necessary: if $n=4k+3$, then $n^2=16k^2+24k+9=16k^2+24k+8+1$, and modulo $8$ this is also $1$ (all of $16k^2,24k,8$ are multiples of $8$) --- so $Q$ can hold with $n\equiv3\pmod4$, where $P$ fails, ruling out A and C. Sufficiency alone rules out D.}
\inv{The true necessary-and-sufficient condition for $Q$ is simply ``$n$ is odd'' --- every odd square is $\equiv1\pmod8$, regardless of which odd residue mod $4$ it comes from.}

%── D4 ──────────────────────────────────────────────────────────────────────────
\item Statement $S$: ``If a quadrilateral has four equal sides, then its diagonals are perpendicular.'' (a) Is $S$ true? (b) State $S$'s converse. Is the converse true?

\ans{(a) True (a rhombus). (b) Converse: ``If a quadrilateral's diagonals are perpendicular, then it has four equal sides.'' The converse is false.}
\method{(a) Let quadrilateral $WXYZ$ have $WX=XY=YZ=ZW$, and let $M$ be the intersection of diagonals $WY$ and $XZ$. Triangles $WXY$ and $WZY$ are congruent (SSS: $WX=WZ$, $XY=ZY$, $WY$ shared), so $\angle XWY=\angle ZWY$, meaning $WY$ bisects $\angle XWZ$; since triangle $WXZ$ is isosceles ($WX=WZ$), this bisector is also the perpendicular bisector of $XZ$. So the diagonals meet at right angles. (b) A kite with $WX=WZ$ but $XY=ZY\neq WX$ (two pairs of adjacent equal sides, not all four equal) has the same SSS congruence argument still apply to triangles $WXY,WZY$ --- so its diagonals are still perpendicular, but it is a counterexample to the converse since its sides are not all equal.}
\inv{``Four equal sides'' is a strictly stronger condition than ``diagonals perpendicular'' --- every rhombus is a kite, but not every kite is a rhombus.}

%── D5 ──────────────────────────────────────────────────────────────────────────
\item Statement $S$: ``If $n$ is a positive integer such that $n^2-1$ is prime, then $n=2$.'' (a) Is $S$ true? Prove or disprove. (b) State $S$'s contrapositive. (c) Is ``$n^2-1$ is prime'' necessary and sufficient for ``$n=2$''?

\ans{(a) True. (b) ``If $n\neq2$, then $n^2-1$ is not prime.'' (c) Yes.}
\method{(a) Factor $n^2-1=(n-1)(n+1)$. For a product of two positive integers to be prime, one factor must equal $1$. Since $n\geq1$, the smaller factor is $n-1$, forcing $n-1=1$, i.e.\ $n=2$; checking, $n=2$ gives $n^2-1=3$, which is indeed prime. For every other $n>2$, both $n-1\geq2$ and $n+1\geq4$, so $(n-1)(n+1)$ is composite. So $n=2$ is true, and it is the unique positive integer. (b) Negate both clauses of $S$ and swap. (c) Since $n=2$ is the unique solution, ``$n^2-1$ prime'' and ``$n=2$'' describe exactly the same set --- necessary and sufficient.}
\inv{Uniqueness turns a one-directional conditional implication into a biconditional equivalence for free.}

\end{enumerate}

%═══════════════════════════════════════════════════════════════════════════════
\section*{Top 5 Patterns Today}
%═══════════════════════════════════════════════════════════════════════════════

\begin{enumerate}[leftmargin=2em, itemsep=4pt]
  \item \textbf{Directional Testing of Necessity vs Sufficiency.}
        ``$P$ is sufficient for $Q$'' means $P \implies Q$; ``$P$ is necessary for $Q$'' means $Q \implies P$. Always evaluate the forward and backward implications independently. \seealso{A1, A3, A5, B1, B5}
  \item \textbf{Strict Inclusions and Counterexample Witnesses.}
        A condition that is strictly more restrictive is sufficient but not necessary (e.g.\ $f'(x)>0$ for strict monotonicity, unit disk in $|x|+|y|\le\sqrt{2}$). Disprove necessity with an edge counterexample. \seealso{A5, A7, B6, B10, D2}
  \item \textbf{Equivalence in Geometry and Algebra.}
        $P \iff Q$ holds when both directions are valid: altitude equality forces equilateral sides, and symmetric quadratic discriminants pin down positive roots. \seealso{A3, A4, A6, A9, C3}
  \item \textbf{Polynomial and Cubic Root Separation.}
        Distinct roots in intervals or cubic extrema requiring $4p^3 > q^2$ demand checking discriminant signs, axis vertices, and interval boundaries together. \seealso{A3, B1, B8, C2, D1}
  \item \textbf{Divisibility and Modular Invariants.}
        $n^2 \equiv 1 \pmod 8 \iff n$ odd; $n^3-n \equiv 0 \pmod{24}$ for odd $n$. Exponent parity in prime factors establishes exact divisibility equivalences. \seealso{A8, A10, B4, B9, D3}
\end{enumerate}

%═══════════════════════════════════════════════════════════════════════════════
\section*{Common Traps to Avoid}
%═══════════════════════════════════════════════════════════════════════════════

\begin{enumerate}[leftmargin=2em, itemsep=4pt]
  \item \textbf{Assuming Sufficiency Automatically Implying Necessity.}
        $f'(x)>0$ guarantees $f$ is strictly increasing, but $f(x)=x^3$ is strictly increasing with $f'(0)=0$; never assume the converse holds without checking. \seealso{A5, B6, B10, D2, D3}
  \item \textbf{Overlooking Sign Flips in Reciprocals and Inequalities.}
        $x > y$ does not imply $1/x < 1/y$ when signs differ ($x=1, y=-2$); sign preservation is essential when inverting inequalities. \seealso{A2, B3, C1, D4}
  \item \textbf{Assuming False Transitivity Across Shared Consequences.}
        $P \implies Q$ and $R \implies Q$ does not imply any logical link between $P$ and $R$ ($x=1$ and $x=-1$ both square to $1$). \seealso{B7, C4}
  \item \textbf{Confusing Contrapositive with Converse.}
        The contrapositive of $a^2+b^2 \equiv 0 \pmod 3$ always shares its truth value; the converse requires its own independent justification. \seealso{B4, B9, C8, D1, D4}
  \item \textbf{Missing Multiplicity or Trivial Degeneracies.}
        In $x^4+x^2+c=0$, $c=0$ yields $x=0$ with multiplicity 2 rather than distinct real roots; strict inequalities at boundaries make or break equivalence. \seealso{A3, B1, B8, C7, D5}
\end{enumerate}

\SpeedClosing{``Speed comes from recognition, not from rushing. Every shortcut is a pattern you have internalised.''}

\end{document}
